\documentclass[11pt]{article}
\usepackage[a4paper,margin=27mm]{geometry}
\usepackage{amsmath,amssymb,mathtools}
\usepackage{fontspec,unicode-math}
\setmainfont{Latin Modern Roman}
\setmathfont{Latin Modern Math}
\usepackage{xurl,hyperref}
\hypersetup{colorlinks=true,urlcolor=blue,linkcolor=blue}
\newcommand{\sref}[2]{\hyperlink{source:#1:#2}{[S#2]}}
\newcommand{\cataloguescope}[1]{\paragraph{Mathematical question.} #1}
\begin{document}
\section{Open Mapping Theorem for Separable Banach Spaces}\label{statement}

\subsection{Definitions and mathematical statement}

\paragraph{Shared notation.}$\mathbb N=\{1,2,\ldots\}$, and $\mathbb Z,\mathbb Q,\mathbb R,\mathbb C$ denote the integers, rationals, reals and complex numbers. Any other conventions are specified locally.


Work over $\mathsf{RCA}_0$, the subsystem of second-order arithmetic with $\Delta^0_1$ comprehension and $\Sigma^0_1$ induction; all structures and orders are countable and coded by sets of integers. $\mathsf{ACA}_0$ adds comprehension for every arithmetical formula (set parameters allowed), and $\mathsf{ATR}_0$ adds arithmetical transfinite recursion along every coded well-order. Reverse-mathematical equivalence means provability of both implications over this base. $\mathsf{WKL}_0$ adds the assertion that every infinite binary tree has an infinite path. A separable Banach space is a complete normed real vector space with a countable dense set. For a bounded linear surjection $T:X\to Y$ between such spaces, the open mapping theorem asserts that $T$ sends every norm-open set to a norm-open set; equivalently there is $\delta>0$ with $B_Y(0,\delta)\subseteq T(B_X(0,1))$.
\paragraph{Statement recorded in the supplied source.}
Is this separable open mapping theorem provable in $\mathsf{RCA}_0$? Is it provable in $\mathsf{WKL}_0$? \sref{statement}{2}
\subsection{Short English statement}
Determine whether either recursive comprehension alone or weak König's lemma suffices for the open mapping theorem on separable Banach spaces.
\subsection{Sources}
\begin{description}
\item[\hypertarget{source:statement:1}{S1}] ULAM.AI and the UnsolvedMath Contributors, UnsolvedMath: A Curated Collection of Open Mathematics Problems, catalogue entry. CC BY 4.0. \url{https://huggingface.co/datasets/ulamai/UnsolvedMath}.
\item[\hypertarget{source:statement:2}{S2}] Stephen G. Simpson, Open Problems in Reverse Mathematics (15 December 1999), Banach Space Theory, first example. \url{https://sgslogic.net/t20/talks/cta/problems/node3.html}.
\end{description}
\label{end:statement}
\end{document}
